\section{Related Work}\label{sec:related}

Convolutional neural networks (CNNs) are the usual starting point for land-use classification of
Sentinel-2 patches. EuroSAT~\cite{helber2019eurosat}, the dataset we trained on, was introduced
together with a benchmark of deep CNNs over its 13 spectral bands and 10 classes, and those CNNs
reached a high overall accuracy. BigEarthNet~\cite{sumbul2019bigearthnet} is a much larger
Sentinel-2 archive in which each patch has several land-cover labels. Its authors found that a
shallow CNN trained on it gave much higher accuracy than a state-of-the-art CNN pre-trained on
ImageNet, a result that suggests that features learned on natural photographs do not fully
transfer to Sentinel-2 images. Because both
benchmarks were introduced with CNN baselines, CNNs are the natural reference when a new model is
tested on Sentinel-2 patches. We used two CNN families. ResNet~\cite{he2016resnet} learns residual
functions, which makes very deep networks easier to optimize. EfficientNet~\cite{tan2019efficientnet}
scales depth, width and input resolution together with one compound coefficient, and its authors
report better accuracy and efficiency than previous CNNs.

Transformers reached image classification later. The Vision Transformer
(ViT)~\cite{dosovitskiy2020vit} applies a pure transformer to a sequence of image patches,
without convolutions. When pre-trained on large amounts of data, it obtained excellent results
compared with state-of-the-art CNNs while needing fewer computational resources to train. Bazi
et~al.~\cite{bazi2021vit} applied ViT to remote-sensing scene classification. In their model,
multihead attention derives long-range relations between pixels, and on four remote-sensing
datasets the model showed promising results compared with state-of-the-art methods. SatMAE~\cite{cong2022satmae}
went one step further and pre-trained transformers for temporal or multispectral satellite
images with a masked autoencoder. Its authors found it beneficial to encode the spectral bands as
groups with their own spectral position encodings, and they report gains in transfer learning,
including land-cover classification. This suggests that a transformer can use extra spectral
bands when its pre-training is designed for them, which is relevant to our 13-band runs.

Self-supervised pre-training addresses a different gap: labels are scarce, while unlabeled
satellite images are plentiful. A review by Wang et~al.~\cite{wang2022sslreview} introduces
self-supervised learning for remote sensing and notes that most of its potential in Earth
observation remains locked. Seasonal Contrast (SeCo)~\cite{manas2021seco} starts from the same
observation that motivates our study: models pre-trained on ImageNet may not generalize to remote
sensing because of the domain gap. SeCo gathers unlabeled images of many locations at different
times and uses invariance to time and position to learn transferable representations; its models outperformed their
ImageNet counterparts on several downstream tasks. SSL4EO-S12~\cite{wang2023ssl4eo} provides a
large, global and multiseasonal corpus of Sentinel-1 and Sentinel-2 images, with pre-trained
models for methods such as MoCo and DINO. We used its ResNet-50 and ViT-S/16 weights as the
self-supervised alternative to ImageNet.

The last line of work concerns what happens when the test region differs from the training
region. Tuia et~al.~\cite{tuia2016domain} note that when training samples come from a different
spatial region than the one being mapped, spectral shifts between the two distributions are
likely to make the model fail. They attribute such shifts to differences in acquisition and
atmospheric conditions or to changes in the observed objects, and they review domain adaptation
methods designed to be robust to such shifts. WILDS~\cite{koh2021wilds} collects real distribution shifts,
including shifts across time and location in satellite imaging, and shows that standard training
gives much lower performance out of distribution than in distribution. Measuring such a shift
also requires an honest test set. Roberts et~al.~\cite{roberts2017cv} show that random
cross-validation on spatially structured data can seriously underestimate the prediction error,
and they recommend block cross-validation, where the data are split by blocks instead of at
random. These studies leave a practical question open for our setting: how two CNNs and a ViT,
trained under one shared protocol, behave when they move from European Sentinel-2 patches to a
semi-arid region of Mexico, measured on hand-labeled patches split by spatial blocks.
